\documentclass[a4paper]{article}

% Math packages
\usepackage[usenames]{color}
\usepackage{xcolor}
\usepackage{forest}
\usepackage{ifxetex,ifluatex,amssymb,amsmath,mathrsfs,amsthm,witharrows,mathtools,mathdots,centernot,stmaryrd}
\usepackage{cancel,bm}
\WithArrowsOptions{displaystyle}
\renewcommand{\qedsymbol}{$\blacksquare$} % end proofs with \blacksquare. Overwrites the defualts.

% Design
\RequirePackage[labelfont=bf]{caption}
\RequirePackage[margin=0.8in]{geometry}
\RequirePackage{multicol}
\RequirePackage[skip=4pt, indent=0pt]{parskip}
\RequirePackage[normalem]{ulem}
\forestset{default}
\renewcommand\labelitemi{$\bullet$}
\RequirePackage{titlesec}
\titleformat{\section}[block]
	{\huge}
	{\dotfill (\textbf{\thesection})\dotfill}
	{0em}
	{}
\usepackage[hidelinks]{hyperref}

\usepackage{../../../tex/packages/mathShortcuts}
%
\makeatletter
\newcommand{\skipitems}[1]{
	\addtocounter{\@enumctr}{#1}
}
\makeatother

\newcommand \middleText[1] {\vfil {\hfil {#1}} \hfil \vfil \newpage}
\newcommand \envendproof{\vspace{-17pt}}
\newcommand \mathenvendproof{\par\vspace{-24pt}}

\newcommand\nsubsection[1]{\subsection*{#1}\addcontentsline{toc}{subsection}{#1}}

\usepackage{titling}
\author{Shahar Perets}
\title{Partial Solutions To Exercises In Group Theory}
\date{\today}
\begin{document}
%	\maketitle
	\huge {\huge\,\\ \vfil {\begin{center}
				\thetitle \\
				\vspace{10pt}
				\large \theauthor \\
				\vspace{3pt}
				\large \thedate
		\end{center}} \vfil	}
	\normalsize
	
	\pagebreak
	
	\section*{Intro}
	Solutions the HW1-10 from 24B in group theory. Each section represent a HW. 


	\tableofcontents
	
	
	\section{}
	\begin{enumerate}
		\item 
		\begin{enumerate}
			\item 
			\item $G \times H$ is Abelian iff $G, H$ are both Abelian\begin{proof}
				Indeed: 
				\begin{itemize}
					\item[$\implies$]If $G \times H$ is Abelian, so $\pi_1, \pi_2$ are group homomorphisms (which preserve group properties like being Abelian) therefore its images $G, H$ are also Abelian. 
					\item[$\impliedby$]If $G, H$ are Abelian, and let $(g_x, h_x) = x, (g_y, h_y) = y \in G \times H$, then: 
					\[ \begin{WithArrows}
						x \cdot y &= (g_x, h_x) \cdot (g_y, h_y) = (g_x \cdot g_y, h_x \cdot h_y) \Arrow[down]{$G, H$ are Abelian}\\
						&= (g_y\cdot g_x, h_y \cdot h_x)  \\
						&= (g_y, h_y) \cdot (g_x, h_x) \\
						&= y \cdot x
					\end{WithArrows} \]
					As intended. 
				\end{itemize}
			\end{proof}
		\end{enumerate}
		\item Assume $\forall g \in G \co g^{2} = e$. We'll prove $G$ is Abelian. \begin{proof}
			Let $x, y \in G$. Clearly $x = x\op \land y = y\op$ (since $y^{2} = x^{2} = 1$). Then $(xy)^{2} = e$ therefore $xy = (xy)\op = y\op x\op = y x$ as intended. 
		\end{proof}
		\item Consider the group $G = \gL_2(\R)$ (the linear space of dimension two with addition) and the groups $H_1 = \Sp((1, 0)), H_2 = \Sp((0, 1))$. Clearly $H_1, H_2$ are subgroups (since they're sub-vector-spaces), though $H_1 \cup H_2$ isn't not a group since $(1, 0) + (0, 1) \notin H_1 \cup H_2$ (addition not closed). 
		\item Consider $\gSL_2(\Z)$. We'll show its a subgroup of $\gGL_2(\R)$. \begin{proof}
			Clearly $\gSL_2(\Z) \subset \gL_2(\R)$. Though since if $\det A = 1$ then $A$ is invertible, we know that $\gSL_2(\Z) \subset \gGL_2(\R)$. It remains to show that $\gSL_2(\Z)$ is a group. It's clear since if $A, B \in \gSL_2(\Z)$ then $\det AB = \det A \det B = 1$ and the coefficients of $AB$ must lie in $\Z$ (matmul definition) therefore it's closed another multiplication. 
			Similarly $\det(A\op) = (\det A)\op = 1$ and $A\op = \frac{\adj A}{\det A} = \adj A$, and clearly the determinant of the minors of $A$ are integers (an integer polynom that gets an integer value yields and integer result) therefore $A\op = \adj A \in \gSL_2(\Z)$ too. 
			
			We already know $\gGL_2(\R)$ multiplication satisfies the conditions for group multiplication, so $\gSL_2(\Z)$ must form a subgroup of $\gGL_2(\R)$. 
		\end{proof}
		\item 
		\item 
	\end{enumerate}
	
	\section{}
	\begin{enumerate}
		\item We'll prove that $g^{n} = 1 \iff \ord(g) \mid n$. \begin{proof}
			If $\ord (g) \mid n$ then $k \cdot \ord(g) = n$ for some $k \in \N_0$. Almost by definition $g^{\ord(g)} = e$ and now:
			\[ g^{n} = g^{\ord (g)\cdot k} = \cl{g^{\ord g}}^{k} = 1^{k} = k \] 
			As intended. 
			
			The other side will be proven via contrapositive. AFSOC $\ord(g)\nmid n \land g^{n} = 1$. Then consider the division with (minimal) remainder $n = k \ord(g) + r$ ($0 < r < \ord (g)$). Then we get that: 
			\[ 1 = g^{n} = g^{k \ord (g) + r} = g^{r} \cdot \cl{g^{\ord (g)}}^{k} = g^{r} \]
			So we found $r$ s.t. $g^{r} = 1$, and $r \neq 0$. Therefore $\ord(g) \le r$, yet $r < \ord (g)$, which yields a contradiction. 
		\end{proof}
		\item Consider $G = \la g \ra$ a cyclic group generated by $g$ with order $\sof G = n$. Let $d \mid n$. Then: 
		\begin{itemize}
			\item $H = \{x \in G \co x^{d} = 1\}$ is a subgroup with order $d$. \begin{proof}
				Clearly $H$ is a subgroup (if $x^{d} = 1 \land y^{d} = 1$ then $(xy)^{d} = x^{d}y^{d} = 1$ and $(x\op)^{d} = (x^{d})\op = 1\op = 1$). 
				We know that $\exists k \in \N_0 \co dk = n$. Now consider $p \cod g^{k}$. Clearly $p^{d} = p^{kd} = p^{n} = 1$. Now consider the subgroup $H = \la p \ra $. It's a cyclic subgroup with order of $d$, since $\ord(p) = d$. We'll show that $d \in H$: $H$ must be cyclic (subgroup of a cyclic group must be cyclic) therefore it's generated by some element $p'$. Note that $p' = p$ otherwise $p'^{z} = p$ for $z > 1$ yet it means that $p^{d - z} = 1$ although $\la p \ra$ has an order of $d$. Hence indeed $H = \la p \ra$, which has the order $d$. 
			\end{proof}
			\item We'll show it's the only subgroup with order $d$. \begin{proof}
				Consider the existence of a (cyclic) group $H'$ with order $d$. Then it's generated by some $p'$ s.t. the minimum argument $\ml$ s.t. $p'^{\ml} = 1 $ is $\ml = d$. This strictly implies $p = g^{k}$ (with the notation $dk = n$ above) as $d \mid n$ (it's given but also known more generally by Lagrange). 
%				TODO
			\end{proof}
		\end{itemize}
		\item 
		\item We'll prove $(\Q, +)$ is not finitely generated. \begin{proof}
			If $(\Q, +)$ was finitely generated by some $g_1 \dots g_n$ where $g_i = \frac{p_i}{q_i}$ ($\gcd(p_i, q_i) = 1$) then the number $g = \frac{1}{1 + \sum_{i = 1}^{n} q_i} \in \Q$ could have not been reached since the denominator of $g$ must be a multiplication of the $q_i$-s (which of course it isn't). We found $g \in \Q$ but not in the group generated by $g_1 \dots g_n$, and via deductio ad absurdum we finished. 
		\end{proof}
		\item We'll prove the double-costs are a partition of $G$. \begin{proof}
			content
		\end{proof}
	\end{enumerate}
	
	\section{}
	\begin{enumerate}
		\item We need to calculate $3^{199} \bmod 101$. Notice that $3^{101} = 3 \mod 101$ therefore in $\Z/101\Z$ we know that $3^{202} \equiv_{101} 3^{2} \equiv 9$. We aim to find $x$ s.t. $x \cdot 3^{3} \equiv = 101k + 9$. Going through the options: 
		\begin{itemize}
			\item For $k = 1$ we get $110 \nmid 3^{3}$
			\item For $k = 2$ we get $221 \nmid 3^{3}$
			\item ...
			\item For $k = 9$ we get $918 \mid 3^{3}$
		\end{itemize}
		Therefore $x \cdot 3^{3} \equiv 918 \implies x \equiv \frac{918}{3^{3}} = 34$. 
		\item Let $G$ be a group, and $H, K$ subgroups. We'll prove that $[H : H \cap K] \le [G : K]$. \begin{proof}
			Consider the function $\phi \co H/(H \cap K) \to G/K$ defined as $\phi(kH) = kG$ (where $k \in H \cap K$). It's clearly independent of choice of representative $k$ since if $(H \cap K)k_1 = (H \cap K)k_2$ then $Hk_1 = Hk_2$ (since $k_1 k_2\op \in H \cap K \subset H$ in this case). Moreover, $\phi$ is injective since if $Hk_1 = Hk_2$ then $(H \cap K)k_1 = (H \cap K)k_2$
			
%			TODO
%			By applying Lagrange's Theorem twice we obtain: 
%			\begin{align*}
%				{\sof H} = \frac{\sof G}{[G : H]} && {\sof{H \cap K}} = \frac{\sof{H}}{[H : H \cap K]}
%			\end{align*}
%			Since $H \le G$ and $H \cap K \le H$ we know that $\sof H \le \sof G \land \sof {H \cap K} \le \sof H$. By combining these inequalities with the mentioned above we get: 
%			\begin{align*}
%				\frac{\sof G}{[G : H]} = \sof H &\ge \sof{H \cap K} = \frac{\sof{H}}{[H : H \cap K]} \\
%				\sof {G} [G : H]&\ge [H : H \cap K] \sof H
%			\end{align*}
%			therefore $\frac{\sof H}{\sof{H \cap K}} \ge 1$. By combining these inequalities with the mentioned above we get: 
%			\begin{align*}
%				[H : H \cap K]\sof{H \cap K} = \sof H &\le \sof G = [G : H] \sof H \\
%				[H : H \cap K] &\le [G : H] \cdot \frac{\sof H}{\sof{H \cap K}} \le [G : H]
%			\end{align*}
		\end{proof}
		\item Let $G$ be a group generated by $S \subset G$. Let $\vphi, \psi \co G \to H$ be homomorphisms s.t. $\forall s \in S \co \phi(s) = \psi(s)$. Then $\phi = \psi$. \begin{proof}
			Let $g \in G$. By definition, exists $s_1 \dots s_n \in S$ s.t. $s_1 + \cdots + s_n = g$. Therefore: 
			\[ \forall g \in G \co \phi(g) = \phi\cl{\sum_{i = 1}^{n}s_i} = \sum_{i = 1}^{n}\phi(s_i) = \sumni \psi(s_i) = \psi\cl{\sumni s_i} = \psi(g) \]\envendproof
		\end{proof}
		\item We'll find all of the homomorphisms $f \co G \to H$ in the following cases: 
		\begin{enumerate}
			\item Let $G = H = \Z$. Claim: $f(z) = az$ for some $a \in \Z$ ($f$ is the quotient map that maps $\Z$ into the normal subgroup $a\Z$). Clearly $f$ is an homomorphism, we'll proof it's the only homomorphism. Indeed, considering some $f \co \Z \to \Z$ homomorphism, denote $f(1) = a$. Then for each $z \in \Z$, because $\Z$ is finitely generated by $1$ ($z = 1 + 1 + \cdots + 1 \doc 1 \cdot z$) we know that $f(1 \cdot z) = z \cdot f(1)$ by definition. Now $f(z) = az$ as intended. 
			
			[btw, $f$ is a ring isomorphism too]
			\item Let $G = \Z, H = \Z/n\Z$. Claim: $f(z) = [anz]$ for $a \in \Z$ are the only homomorphisms $\Z \to \Z/n\Z$. It's clearly a homomorphisms. And the proof the existence is similar to the one in (A)
			\item 
			\item 
		\end{enumerate}
		\item Proof/Disproof the following claims: 
		\begin{itemize}
			\item Claim: $(\Z, +) \cong (\Q, +)$. \begin{proof}[Disproof]
				Consider the set of generators of $\Z$, which is $S = \{1_\Z\}$. Assume for the sake of contradiction that there exists and isomorphism $f \co \Z\to \Q$. Therefore $f(S)$ is the set of generators of $\Q$. Identities are send to identities by isomorphisms and thereafter $f(S) = \{1_\Q\}$. But $\frac{1}{2} \in \Q$ is clearly not generated by $1_\Q$, and ad absurdum the claim is false. 
			\end{proof}
			\item Claim: $(\Q, +) \cong (\Q^\times, \cdot)$. \begin{proof}[Disproof]
				In $\Q^\times$ there exists a root os unity $-1$ (since $-1 \cdot -1 = 1$). Therefore the isomorphism $f \co \Q \to \Q^\times$ preserve the root of unity $f(-1) \cdot f(-1) = 1_\Q = 0$. But in $\Q$ there are not root of unity (if $a + a = 0$ then $2a = 0$ and $a = 0$) therefore $f(-1) = 0$, meanwhile $f(1) = f(1_{\Q^{\times}}) = 1_\Q = 0$ therefore $f$ is not injective, which is a contradiction. 
			\end{proof}
		\end{itemize}
		\item We'll prove that up to an isomorphism, there's only one group with the order of $p$ ($p$ is prime). \begin{proof}
			It's sufficient to prove that every group with a prime order must be cyclic, since all of the finite cyclic groups isomorphic to $\Z/p\Z$. 
			
			Let $G$ be a finite group with the prime order of $p$. Consider the cyclic subgroup $H \le G$ generated by $a \neq 1_G$ ($G$ is nontrivial since $p \neq 1$ therefore such $a$ must exists). By Lagrange $\ord G = \ord H \cdot [G : H]$. But $\ord G = p$ is prime, therefore either $\ord H = 1$ or $\ord H = p$. But $a, 1_G \in H$ therefore $\ord H \neq 1$. Finally $\ord H = \ord G = p$, hence $H = G$ and $G$ is cyclically generated by $a$. 
		\end{proof}
	\end{enumerate}
	
	\section{}
	
	
	
	
	
	
	
\end{document}
